\section{结论}
\label{sec:conclusion}

本文介绍了 PhysFormer，一种物理引导的可解释 Transformer，旨在调和 VPP 预测中预测准确性、结构可解释性和物理合规性之间的根本权衡。通过有界物理激活-残差（BPAR）机制，将显式物理映射流与神经 Transformer 融合，PhysFormer 在结构上约束了神经残差的搜索空间。这种方法在不采用不可微启发式裁剪的情况下，有效地消除了非物理异常，实现了绝对的边界合规性（BVR=0.00\%），并通过将必要的结构化边界修正限制在物理上可忽略不计的量级（$2.6 \times 10^{-5}$ MW），保持了时间爬坡的完整性。

至关重要的是，PhysFormer 并没有为了安全而牺牲准确性，而是引入了一个在架构上正交的物理引导因果耦合（PGCC）模块。通过显式的门控响应正则化，PGCC 成功地将高度相关的物理关系对齐到时间门控变量上。这种结构化的对齐产生了优异的全局相关性（$r=0.8306$），并且在活跃的日间发电时段保持了稳健的耦合（$r=0.4058$）。此外，即使在高达 50\% 的分布外传感器噪声下，这种结构相关性仍具有韧性（$r > 0.81$）。详尽的消融实验证实，这种结构透明度的实现并没有蚕食深度不受约束残差跟踪的预测能力（MSE 维持在 0.0128 不变）。

因此，PhysFormer 确立了一个占优的多维帕累托地位。它从根本上克服了困扰纯数据驱动流水线的结构合规性缺陷，确保了在功能上等同于最强不受约束基准模型的预测性能，同时在数学上保证了绝对的运行安全性和透明的物理可追溯性。

未来的工作将把 BPAR 架构扩展到涵盖分布式能源资产的多年度退化轨迹，并探索将这些符合力学规律的预测流形直接嵌入到下游基于强化学习的 VPP 调度和交易算法中 \cite{Ma2025, Arabzadeh2025}。
